\addcontentsline{toc}{subsubsection}{Solution 7. The Harmonic Squeeze}

\begin{dlsummary}[title={\bfseries Solution 7. The Harmonic Squeeze},breakable]

\begin{proof}
Set
\[
g(x)=\frac1{f(x)}.
\]
Since all function values are positive, taking reciprocals in the given
inequalities gives
\[
g(x+y)\ge g(x)+g(y),
\qquad
g(xy)\le g(x)g(y),
\]
and the condition on $f(2)$ gives $g(2)\le2$.

\medskip
\noindent
\textbf{Step 1. Lock the dyadic values.}

Taking $x=y=1$ in submultiplicativity yields
\[
g(1)\le g(1)^2,
\]
so $g(1)\ge1$. Therefore
\[
2\ge g(2)\ge2g(1)\ge2,
\]
which implies $g(1)=1$ and $g(2)=2$. By induction, superadditivity and
submultiplicativity give
\[
g(2^n)=2^n
\qquad(n=0,1,2,\ldots). \tag{1}
\]
Moreover, $g$ is strictly increasing, since
\[
g(y+x)\ge g(y)+g(x)>g(y)
\qquad(x,y>0).
\]

\medskip
\noindent
\textbf{Step 2. Establish an upper bound.}

Fix $t>0$. For sufficiently large $n$, let
\[
m_n=\left\lfloor\frac{2^n}{t}\right\rfloor.
\]
Since $m_nt\le2^n$, monotonicity, superadditivity, and (1) yield
\[
2^n=g(2^n)\ge g(m_nt)\ge m_ng(t).
\]
As $2^n/m_n\to t$, it follows that
\[
g(t)\le t\qquad(t>0). \tag{2}
\]

\medskip

For every positive integer $k$, we have $g(k)=k$. If $x>1$, repeated
submultiplicativity and monotonicity give
\[
g(x)^n\ge g(x^n)
\ge g\bigl(\lfloor x^n\rfloor\bigr)
=\lfloor x^n\rfloor.
\]
Taking $n$-th roots and letting $n\to\infty$ gives $g(x)\ge x$.
Together with (2), this proves $g(x)=x$ for every $x>1$.

For $0<x<1$, submultiplicativity gives
\[
1=g(1)\le g(x)g\left(\frac1x\right)=\frac{g(x)}x,
\]
where the last equality follows from the result for $1/x>1$. Thus
$g(x)\ge x$, and (2) forces equality. Therefore $g(x)=x$ for all
$x>0$, and hence $f(x)=1/x$ for all $x>0$.
\end{proof}

\begin{proof}[Verification]
For $f(x)=1/x$,
\[
f(x+y)\bigl(f(x)+f(y)\bigr)
=\frac1{xy}
=f(x)f(y)
\le \frac1{xy}=f(xy),
\]
and $f(2)=1/2$. Thus the function satisfies all the conditions and is
the unique solution.
\end{proof}

\end{dlsummary}